% REPLACES: main.tex lines 2585-2636 ("\subsection{A.8 Numerical parameter declarations} ..." to "... every number reported in the text.")
% Numbers in this fragment: \NUM{atomless_epsilon_C}.
% The cost entry leaves each vector: the preparation cost is 0 with probability \rho
% and c_H otherwise. If the integrator instead keeps a cost slot in each tuple for
% registry compatibility, its value is 0 in all three vectors. Every other declared
% value is copied unchanged from the current text.
\subsection{A.8 Numerical parameter declarations}\label{pa-parameters}

In each vector the preparation cost is \(0\) with probability \(\rho\) and
\(c_H\) otherwise.

The benchmark parameter vector is

\[
\begin{aligned}
(h,\ell,p,\rho,c_H,b,k)
&=(10,1,0.5,0.25,6,2,0.02),\\
(r_0,r_1,r_2)&=(1.2,3,3.6).
\end{aligned}
\tag{A.55}\label{eq:paper-a-55}
\]

The moderate-value vector is

\[
\begin{aligned}
(h,\ell,p,\rho,c_H,b,k)
&=(2,1,0.5,0.25,0.89,2,0.002),\\
(r_0,r_1)&=(1.05,1.5).
\end{aligned}
\tag{A.56}\label{eq:paper-a-56}
\]

For complementary signals I use

\[
\begin{aligned}
(h,\ell,p,\rho,c_H,b,k)
&=(10,1,0.5,0.85,7.14,2,0.015),\\
(r_0,r_1,a,d)
&=(1.1,2.3,0.70,0.75).
\end{aligned}
\tag{A.57}\label{eq:paper-a-57}
\]

The atomless-cost parameter is \(\varepsilon_C=\NUM{atomless_epsilon_C}\).
It measures the low component's support from zero and the high component's
support around \(c_H\), as in \eqref{eq:paper-a-24}.
The atomless-value half-width is \(\varepsilon_V=\) 0.05. These are separate
experiments.

The exact reserve events are
\(p_L=h\) at \(r=r_0\) and \(p_H=\sqrt{2r_1(h-c_H/M)-r_1^2}\) at
\(r=r_1\), evaluated from these identities under the benchmark vector,
with the tie rule of Section 2. Online Appendix C defines the exact
input values, the derived quantities, their formatting, and the
acceptance criteria behind every number reported in the text.
